% LTeX: enabled=false
% Diapositivas de ejemplo: una por cada patrón de la plantilla (Plantilla_beamer.md, sección 5).
% Sirven de molde para el contenido real y se borran cuando ya no hagan falta.

\seccion{Patrones de la plantilla}

% ===============================================================
% Texto y figura
% ===============================================================
\begin{frame}{Texto y figura}
    \begin{columns}[T]
        \begin{column}{0.5\textwidth}
            \begin{itemize}
                \item Punto principal de la diapositiva.
                \item Un segundo punto, breve.
                \item Un tercero, si hace falta.
            \end{itemize}
        \end{column}
        \begin{column}{0.45\textwidth}
            \centering
            \includegraphics[width=0.7\linewidth]{logo_keralytix.png}
        \end{column}
    \end{columns}
\end{frame}

% ===============================================================
% Diagrama
% ===============================================================
\begin{frame}{Diagrama}
    Una línea que presenta el diagrama.
    \vspace{0.6cm}
    \begin{center}
        \begin{tikzpicture}[node distance=1cm]
            \node[caja] (a) {Mapas corneales y variables clínicas};
            \node[caja, right=of a, fill=azul!15, very thick] (b) {Modelo multimodal};
            \node[caja, right=of b, draw=verde] (c) {Informe con mapa de calor};
            \draw[flecha] (a) -- (b);
            \draw[flecha] (b) -- (c);
        \end{tikzpicture}
    \end{center}
    \idea{Nota o idea clave al pie del diagrama.}
\end{frame}

% ===============================================================
% Tabla
% ===============================================================
\begin{frame}{Tabla}
    \begin{tabla}{colspec = {Q[l,m,4cm] X[l,m]}, rows = {font=\small}}
        Columna & Descripción \\
        Primera fila & Texto de ejemplo. \\
        Segunda fila & Texto de ejemplo. \\
        Tercera fila & Texto de ejemplo. \\
    \end{tabla}
\end{frame}

% ===============================================================
% Bloques
% ===============================================================
\begin{frame}{Bloques}
    \begin{columns}[T]
        \begin{column}{0.47\textwidth}
            \begin{block}{Primer bloque}
                Contenido del primer bloque.
            \end{block}
        \end{column}
        \begin{column}{0.47\textwidth}
            \begin{block}{Segundo bloque}
                Contenido del segundo bloque.
            \end{block}
        \end{column}
    \end{columns}
\end{frame}

% ===============================================================
% Lista
% ===============================================================
\begin{frame}{Lista}
    \begin{itemize}
        \item Primer punto.
        \item Segundo punto.
        \begin{itemize}
            \item Un subpunto.
        \end{itemize}
        \item Tercer punto.
    \end{itemize}
    \vspace{0.4cm}
    \pendiente{Así se marca lo que falta completar.}
\end{frame}
